\documentclass[11pt,a4paper]{article}

% ===================== POLICES & LANGUE =====================
\usepackage{fontspec}
\setmainfont{Latin Modern Roman}
\setsansfont{Latin Modern Sans}
\setmonofont{Latin Modern Mono}
\usepackage{polyglossia}
\setdefaultlanguage{french}

% ===================== MISE EN PAGE =====================
\usepackage[margin=2.5cm]{geometry}
\usepackage{amsmath,amssymb,graphicx,xcolor,booktabs,hyperref,tcolorbox,enumitem,fancyhdr,array,pifont,longtable,multirow,listings}
\tcbuselibrary{theorems,breakable,skins}
\usepackage{tikz}
\usetikzlibrary{shapes.geometric,arrows.meta,positioning,calc,backgrounds,fit}

\newcommand{\cmark}{\ding{51}}
\newcommand{\xmark}{\ding{55}}
\hypersetup{colorlinks=true,linkcolor=blue!60!black,urlcolor=blue!60!black,citecolor=blue!60!black}

% ===================== BOÎTES =====================
\newtcolorbox{objectifbox}[1][]{colback=blue!5,colframe=blue!50,fonttitle=\bfseries,title={#1},breakable}
\newtcolorbox{infobox}[1][]{colback=green!5,colframe=green!50!black,fonttitle=\bfseries,title={#1},breakable}
\newtcolorbox{attentionbox}[1][]{colback=orange!5,colframe=orange!70,fonttitle=\bfseries,title={#1},breakable}
\newtcolorbox{retenirbox}[1][]{colback=yellow!10,colframe=yellow!50!black,fonttitle=\bfseries,title={#1},breakable}
\newtcolorbox{prereqbox}[1][]{colback=cyan!5,colframe=cyan!50!black,fonttitle=\bfseries,title={#1},breakable}
\newtcolorbox{projetbox}[1][]{colback=teal!5,colframe=teal!50!black,fonttitle=\bfseries,title={#1},breakable}

% ===================== CODE (listings) =====================
\definecolor{codebg}{gray}{0.96}
\definecolor{kwcolor}{RGB}{0,90,160}
\definecolor{strcolor}{RGB}{160,55,30}
\definecolor{comcolor}{RGB}{95,120,95}
\definecolor{ndcolor}{RGB}{120,80,150}

\lstdefinelanguage{JS}{
  morekeywords={let,const,var,function,return,if,else,for,of,in,while,do,true,false,null,undefined,typeof,break,continue,new,this},
  morekeywords=[2]{console,log,table,document,Math,window,alert,prompt,forEach,map,filter,push,pop,querySelector,querySelectorAll,addEventListener,classList,textContent,innerHTML,createElement,appendChild,getElementById,includes},
  sensitive=true,
  morecomment=[l]{//},
  morecomment=[s]{/*}{*/},
  morestring=[b]",
  morestring=[b]',
  morestring=[b]`,
}

\lstdefinestyle{js}{
  language=JS,
  basicstyle=\ttfamily\small,
  keywordstyle=\color{kwcolor},
  keywordstyle=[2]\color{ndcolor},
  commentstyle=\color{comcolor}\itshape,
  stringstyle=\color{strcolor},
  backgroundcolor=\color{codebg},
  frame=single,
  rulecolor=\color{gray!45},
  breaklines=true,
  breakatwhitespace=false,
  showstringspaces=false,
  numbers=none,
  tabsize=2,
  columns=fullflexible,
  keepspaces=true,
  literate=%
    {é}{{é}}1 {è}{{è}}1 {à}{{à}}1 {â}{{â}}1 {ê}{{ê}}1
    {î}{{î}}1 {ô}{{ô}}1 {û}{{û}}1 {ù}{{ù}}1 {ç}{{ç}}1
    {ï}{{ï}}1 {ë}{{ë}}1 {É}{{É}}1 {È}{{È}}1 {À}{{À}}1
    {œ}{{œ}}1 {°}{{°}}1 {²}{{²}}1
}

\lstdefinestyle{html}{
  basicstyle=\ttfamily\small,
  backgroundcolor=\color{codebg},
  frame=single,
  rulecolor=\color{gray!45},
  breaklines=true,
  showstringspaces=false,
  numbers=none,
  tabsize=2,
  columns=fullflexible,
  keepspaces=true,
  literate=%
    {é}{{é}}1 {è}{{è}}1 {à}{{à}}1 {ê}{{ê}}1 {ç}{{ç}}1 {î}{{î}}1
}

\lstset{style=js}
\newcommand{\code}[1]{\texttt{#1}}

% ===================== EN-TÊTES =====================
\pagestyle{fancy}
\fancyhf{}
\fancyhead[L]{\small Jàng Afrig --- Lab 7 : DOM et événements}
\fancyhead[R]{\small UCAD --- L1 DUT1-Info}
\fancyfoot[C]{\thepage}
\fancyfoot[L]{\small Dr. El Hadji Bassirou TOURÉ}
\fancyfoot[R]{\small 2025--2026}
\renewcommand{\headrulewidth}{0.4pt}
\renewcommand{\footrulewidth}{0.4pt}

% ===================== TITRE =====================
\title{%
\vspace{-1cm}
\textbf{\LARGE Fondamentaux du Web}\\
\textbf{\LARGE HTML, CSS, JavaScript}\\[0.8cm]
\textbf{\huge Lab 7 --- DOM et événements}\\[0.3cm]
\textbf{\Large Le premier quiz Jàng Afrig jouable}\\[0.5cm]
\large Licence 1 DUT1-Informatique --- 2025--2026}
\date{}
\author{}

\begin{document}
\maketitle
\thispagestyle{empty}

\begin{prereqbox}[Informations pratiques]
\textbf{Le grand jour.} Jusqu'ici, tout votre code s'affichait dans la console. Aujourd'hui, JavaScript va enfin \textbf{modifier la page visible} et \textbf{réagir aux clics}. À la fin de la séance, Jàng Afrig sera un \emph{vrai quiz jouable} du début à la fin.

\smallskip
\textbf{Avant de commencer :}
\begin{itemize}[leftmargin=*,noitemsep]
  \item \code{data.js} (Lab 6, tag \code{v0.6}) avec les pays et \code{choisirAuHasard}.
  \item \code{script.js} branché, et vos pages HTML/CSS des Labs 1 à 4.
  \item Les méthodes de tableaux du Lab 6 (\code{forEach}, accès par index).
\end{itemize}

\smallskip
\textbf{Validation de fin :} on lance le quiz, on répond à 10 questions, le score s'affiche, on peut rejouer. Tag \code{v0.7} poussé + rapport déposé.
\end{prereqbox}

\begin{objectifbox}[Objectifs du Lab 7]
À la fin, vous saurez :
\begin{enumerate}[leftmargin=*,noitemsep]
  \item Comprendre le \textbf{DOM} : la page vue comme un arbre par JavaScript.
  \item \textbf{Sélectionner} un élément (\code{querySelector}, \code{querySelectorAll}).
  \item \textbf{Lire et modifier} le contenu d'un élément (\code{textContent}, \code{value}).
  \item Afficher/masquer des éléments avec \code{classList}.
  \item \textbf{Réagir} aux clics avec \code{addEventListener}.
  \item Assembler tout cela pour faire un \textbf{quiz Capitales} complet et jouable.
\end{enumerate}
\end{objectifbox}

\tableofcontents
\bigskip\hrule\bigskip

\section*{Rappel express des Labs 5 et 6}
\addcontentsline{toc}{section}{Rappel express des Labs 5 et 6}
Vous savez écrire des fonctions, décider avec \code{if}, parcourir avec \code{forEach}, et vous avez \code{data.js} avec un tableau de pays et la fonction \code{choisirAuHasard}. On va brancher tout ça sur la page.

% ===================== SECTION 1 =====================
\section{Le DOM : la page vue comme un arbre}

Quand le navigateur lit votre HTML, il en construit une représentation en mémoire, organisée comme un \textbf{arbre}. On l'appelle le \textbf{DOM} (\emph{Document Object Model}). \emph{Analogie :} l'arbre généalogique d'une famille --- une balise \code{<body>} a des « enfants » (\code{<header>}, \code{<main>}\ldots), qui ont eux-mêmes des enfants.

\begin{center}
\begin{tikzpicture}[
  font=\footnotesize,
  every node/.style={rectangle, rounded corners=2pt, draw, thick, inner sep=3pt},
  html/.style={fill=blue!8, draw=blue!55},
  branche/.style={fill=teal!8, draw=teal!55!black},
  feuille/.style={fill=gray!6, draw=gray!55},
  ed/.style={draw, gray!60!black, thick}
]
  \node[html]    (html)   at (0,3.6)      {\code{<html>}};
  \node[branche] (head)   at (-3.4,2.4)   {\code{<head>}};
  \node[branche] (body)   at (2.4,2.4)    {\code{<body>}};
  \node[feuille] (title)  at (-4.7,1.1)   {\code{<title>}};
  \node[feuille] (script) at (-2.1,1.1)   {\code{<script>}};
  \node[feuille] (header) at (0.5,1.1)    {\code{<header>}};
  \node[branche] (main)   at (4.3,1.1)    {\code{<main>}};
  \node[feuille] (acc)    at (2.6,-0.3)   {\code{\#accueil}};
  \node[feuille] (quiz)   at (4.4,-0.3)   {\code{\#quiz}};
  \node[feuille] (res)    at (6.3,-0.3)   {\code{\#resultat}};

  \draw[ed] (html) -- (head);
  \draw[ed] (html) -- (body);
  \draw[ed] (head) -- (title);
  \draw[ed] (head) -- (script);
  \draw[ed] (body) -- (header);
  \draw[ed] (body) -- (main);
  \draw[ed] (main) -- (acc);
  \draw[ed] (main) -- (quiz);
  \draw[ed] (main) -- (res);
\end{tikzpicture}

\textit{Le DOM : un arbre de balises. JavaScript peut s'y promener pour lire ou modifier n'importe quel élément.}
\end{center}

\begin{infobox}[À voir dans DevTools]
Ouvrez votre page, appuyez sur \textbf{F12}, et allez dans l'onglet \textbf{Elements} (ou \emph{Éléments}). Vous voyez cet arbre, et vous pouvez déplier chaque branche. C'est la page telle que JavaScript la « voit ».
\end{infobox}

% ===================== SECTION 2 =====================
\section{Sélectionner un élément}

Avant de modifier un élément, il faut d'abord le \textbf{désigner}. L'outil principal pour cela est la fonction \code{querySelector} : on lui donne un \textbf{sélecteur CSS}, exactement comme au CM2.

\begin{lstlisting}[style=js]
const titre = document.querySelector("h1");        // par balise
const bouton = document.querySelector("#btn-commencer"); // par id (#)
const premiere = document.querySelector(".reponse");     // 1re classe (.)
\end{lstlisting}

Pour récupérer \emph{tous} les éléments correspondants (par exemple les 4 boutons de réponse), on utilise \code{querySelectorAll}, qui renvoie une \textbf{liste} :

\begin{lstlisting}[style=js]
const boutons = document.querySelectorAll(".reponse"); // les 4 boutons
console.log(boutons.length); // 4
\end{lstlisting}

\begin{retenirbox}[Le pont avec le CSS]
Les sélecteurs sont \emph{exactement} ceux du CSS : \code{h1} (balise), \code{\#identifiant} (id), \code{.classe} (classe). Ce que vous saviez cibler en CSS, vous savez maintenant le saisir en JavaScript.
\end{retenirbox}

% ===================== SECTION 3 =====================
\section{Lire et modifier le contenu}

\subsection{Le texte d'un élément : \texttt{textContent}}
\code{textContent} permet de \textbf{lire} ou de \textbf{changer} le texte d'un élément.

\begin{lstlisting}[style=js]
const titre = document.querySelector("h1");

console.log(titre.textContent);       // lit : "Jàng Afrig"
titre.textContent = "Bienvenue !";    // modifie : la page change aussitôt
\end{lstlisting}

\subsection{La valeur d'un champ : \texttt{value}}
Pour un champ de saisie (\code{<input>}), ce que l'utilisateur a tapé se trouve dans \code{.value}.

\begin{lstlisting}[style=js]
const champ = document.querySelector("#pseudo");
console.log(champ.value); // ce que l'utilisateur a écrit dans le champ
\end{lstlisting}

\begin{infobox}[Démo en direct]
On sélectionne le \code{<h1>} et on tape \code{titre.textContent = "Test"} : le titre change instantanément à l'écran. C'est la première fois que JavaScript \emph{touche} la page !
\end{infobox}

% ===================== SECTION 4 =====================
\section{Afficher et masquer : \texttt{classList}}

JavaScript peut ajouter ou retirer des \textbf{classes CSS} à un élément, en direct, avec \code{classList} :
\begin{itemize}[leftmargin=*,noitemsep]
  \item \code{element.classList.add("cachee")} --- ajoute la classe ;
  \item \code{element.classList.remove("cachee")} --- retire la classe ;
  \item \code{element.classList.toggle("cachee")} --- l'ajoute si absente, la retire si présente.
\end{itemize}

\textbf{L'astuce qui change tout.} Si on crée en CSS une classe qui cache un élément :
\begin{lstlisting}[style=html]
.cachee { display: none; }
\end{lstlisting}
\ldots alors ajouter/retirer cette classe en JavaScript permet de \textbf{montrer ou cacher} n'importe quelle partie de la page. C'est exactement ce qu'on va utiliser pour passer de l'écran d'accueil à l'écran du quiz.

% ===================== SECTION 5 =====================
\section{Créer un élément (pour aller plus loin)}

On peut aussi fabriquer un nouvel élément et l'ajouter à la page :

\begin{lstlisting}[style=js]
const liste = document.querySelector("#mes-villes");

const item = document.createElement("li"); // créer un <li>
item.textContent = "Dakar";                // lui donner du texte
liste.appendChild(item);                   // l'accrocher dans la liste
\end{lstlisting}

\code{createElement} fabrique l'élément, \code{appendChild} l'accroche comme enfant. On n'en aura pas besoin pour le quiz d'aujourd'hui (nos 4 boutons existent déjà en HTML), mais c'est bon à connaître.

% ===================== SECTION 6 =====================
\section{Réagir aux clics : \texttt{addEventListener}}

Voici le cœur de l'interactivité. \code{addEventListener} permet de dire : \emph{« quand cet événement se produit sur cet élément, exécute cette fonction »}. \emph{Analogie :} on installe une sonnette ; quand on appuie (l'événement \code{click}), elle déclenche une action.

\begin{lstlisting}[style=js]
const bouton = document.querySelector("#btn-commencer");

bouton.addEventListener("click", () => {
  console.log("Le bouton a été cliqué !");
});
\end{lstlisting}

Les événements les plus courants : \code{click} (clic), \code{submit} (envoi de formulaire), \code{input} (saisie au clavier), \code{keydown} (touche pressée).

\paragraph{Savoir \emph{quel} élément a été cliqué : \texttt{event.target}.} La fonction reçoit un objet \code{event} qui décrit ce qui s'est passé. \code{event.target} est l'élément précis sur lequel on a cliqué --- très utile quand plusieurs boutons partagent le même comportement.

\begin{lstlisting}[style=js]
const boutons = document.querySelectorAll(".reponse");

boutons.forEach((bouton) => {
  bouton.addEventListener("click", (event) => {
    console.log("Tu as cliqué : " + event.target.textContent);
  });
});
\end{lstlisting}

\begin{infobox}[Démo en direct]
On installe l'écouteur ci-dessus, puis on clique sur chacun des boutons : la console affiche le texte du bouton cliqué. \textbf{Notre quiz fonctionnera exactement sur ce principe.}
\end{infobox}

% ===================== SECTION 7 =====================
\section{On assemble : le premier quiz Jàng Afrig}

Place à la construction. On va transformer Jàng Afrig en un quiz à \textbf{trois écrans} : Accueil, Quiz, Résultat. Pour garder le score sans recharger la page, on rassemble ces trois écrans dans \code{index.html} sous forme de \code{<section>}, et on n'en montre qu'une à la fois (avec la classe \code{cachee}). C'est la technique des applications web modernes : \emph{une seule page, plusieurs vues}.

\begin{center}
\begin{tikzpicture}[
  font=\footnotesize,
  etat/.style={rectangle, rounded corners=4pt, draw, very thick, align=center, minimum width=2.4cm, minimum height=1.2cm},
  arr/.style={-{Stealth[length=3mm]}, very thick, gray!70!black},
  lab/.style={font=\scriptsize, text=gray!45!black, align=center}
]
  \node[etat, draw=blue!55, fill=blue!6] (acc) {\textbf{Accueil}\\\scriptsize saisir pseudo};
  \node[etat, draw=teal!55!black, fill=teal!6, right=3.2cm of acc] (quiz) {\textbf{Quiz}\\\scriptsize question + 4 choix};
  \node[etat, draw=orange!70, fill=orange!8, right=3.2cm of quiz] (res) {\textbf{Résultat}\\\scriptsize score final};

  \draw[arr] (acc) -- node[lab,above]{clic\\« Mode Capitales »} (quiz);
  \draw[arr] (quiz) -- node[lab,above]{10\textsuperscript{e}\\réponse} (res);
  \draw[arr] (quiz) to[out=120,in=60,looseness=5] node[lab,above]{réponse\\suivante} (quiz);
  \draw[arr] (res) to[out=-120,in=-60] node[lab,below]{clic « Rejouer »} (acc);
\end{tikzpicture}

\textit{La navigation du quiz : trois écrans, pilotés par les clics.}
\end{center}

\subsection{Étape 1 --- restructurer \texttt{index.html} en trois écrans}

Remplacez le contenu du \code{<body>} de \code{index.html} par ceci :

\begin{lstlisting}[style=html]
<body>
  <header><h1>Jàng Afrig</h1></header>

  <main>
    <!-- ÉCRAN 1 : ACCUEIL -->
    <section id="accueil">
      <p>Teste tes connaissances sur les capitales africaines !</p>
      <label for="pseudo">Ton pseudo :</label>
      <input type="text" id="pseudo" placeholder="Ex : Aminata">
      <button id="btn-commencer">Mode Capitales</button>
    </section>

    <!-- ÉCRAN 2 : QUIZ -->
    <section id="quiz" class="cachee">
      <p id="progression">Question 1 / 10</p>
      <h2 id="enonce">Énoncé de la question</h2>
      <div id="reponses">
        <button class="reponse">Réponse A</button>
        <button class="reponse">Réponse B</button>
        <button class="reponse">Réponse C</button>
        <button class="reponse">Réponse D</button>
      </div>
    </section>

    <!-- ÉCRAN 3 : RÉSULTAT -->
    <section id="resultat" class="cachee">
      <h2>Quiz terminé !</h2>
      <p id="score-final">Ton score</p>
      <button id="btn-rejouer">Rejouer</button>
    </section>
  </main>
</body>
\end{lstlisting}

Puis ajoutez cette ligne à la fin de \code{style.css} :
\begin{lstlisting}[style=html]
.cachee { display: none; }
\end{lstlisting}

\subsection{Étape 2 --- compléter \texttt{data.js} (10 pays + mélange)}

Pour un quiz de 10 questions, étoffez le tableau \code{pays} à dix pays, et ajoutez une fonction \code{melanger} (pour brouiller l'ordre des propositions) :

\begin{lstlisting}[style=js]
const pays = [
  { cca3: "SEN", nom: "Sénégal",        capitale: "Dakar" },
  { cca3: "MLI", nom: "Mali",           capitale: "Bamako" },
  { cca3: "GHA", nom: "Ghana",          capitale: "Accra" },
  { cca3: "KEN", nom: "Kenya",          capitale: "Nairobi" },
  { cca3: "MAR", nom: "Maroc",          capitale: "Rabat" },
  { cca3: "NGA", nom: "Nigéria",        capitale: "Abuja" },
  { cca3: "EGY", nom: "Égypte",         capitale: "Le Caire" },
  { cca3: "CIV", nom: "Côte d'Ivoire",  capitale: "Yamoussoukro" },
  { cca3: "TUN", nom: "Tunisie",        capitale: "Tunis" },
  { cca3: "ETH", nom: "Éthiopie",       capitale: "Addis-Abeba" }
];

function choisirAuHasard(tableau) {
  const index = Math.floor(Math.random() * tableau.length);
  return tableau[index];
}

// Mélange un tableau (algorithme classique d'échange)
function melanger(tableau) {
  const copie = [...tableau]; // copie du tableau (à prendre tel quel, CM3)
  for (let i = copie.length - 1; i > 0; i--) {
    const j = Math.floor(Math.random() * (i + 1));
    const temp = copie[i];
    copie[i] = copie[j];
    copie[j] = temp;
  }
  return copie;
}
\end{lstlisting}

\subsection{Étape 3 --- générer les questions (dans \texttt{script.js})}

Une question = un pays « bonne réponse », sa capitale, et trois autres capitales en distracteurs, le tout mélangé.

\begin{lstlisting}[style=js]
function genererQuestion() {
  const bonPays = choisirAuHasard(pays);
  const bonneReponse = bonPays.capitale;

  // on part de la bonne réponse, puis on ajoute 3 distracteurs distincts
  const propositions = [bonneReponse];
  while (propositions.length < 4) {
    const autre = choisirAuHasard(pays);
    if (!propositions.includes(autre.capitale)) {
      propositions.push(autre.capitale);
    }
  }

  return {
    enonce: `Quelle est la capitale du pays : ${bonPays.nom} ?`,
    bonneReponse: bonneReponse,
    propositions: melanger(propositions)
  };
}

function genererQuestions(nombre) {
  const liste = [];
  for (let i = 0; i < nombre; i++) {
    liste.push(genererQuestion());
  }
  return liste;
}
\end{lstlisting}

\begin{infobox}[Deux petites nouveautés, à prendre telles quelles]
\code{while (condition) \{ ... \}} signifie \emph{« tant que la condition est vraie, répète »} (vous l'avez vu en pseudo-code). \code{tableau.includes(x)} renvoie \code{true} si le tableau contient déjà \code{x} : on s'en sert pour éviter deux fois la même capitale.
\end{infobox}

\subsection{Étape 4 --- la logique des écrans (\texttt{script.js})}

\begin{lstlisting}[style=js]
// --- On saisit les éléments de la page (une seule fois) ---
const ecranAccueil  = document.querySelector("#accueil");
const ecranQuiz     = document.querySelector("#quiz");
const ecranResultat = document.querySelector("#resultat");

const champPseudo  = document.querySelector("#pseudo");
const btnCommencer = document.querySelector("#btn-commencer");
const btnRejouer   = document.querySelector("#btn-rejouer");

const elProgression  = document.querySelector("#progression");
const elEnonce       = document.querySelector("#enonce");
const elScoreFinal   = document.querySelector("#score-final");
const boutonsReponse = document.querySelectorAll(".reponse"); // 4 boutons

// --- L'état du jeu ---
let questions = [];
let indexQuestion = 0;
let score = 0;
let pseudo = "";

// --- N'afficher qu'un seul écran ---
function afficherEcran(ecran) {
  ecranAccueil.classList.add("cachee");
  ecranQuiz.classList.add("cachee");
  ecranResultat.classList.add("cachee");
  ecran.classList.remove("cachee");
}

// --- Afficher la question courante ---
function afficherQuestion() {
  const q = questions[indexQuestion];
  elProgression.textContent = `Question ${indexQuestion + 1} / ${questions.length}`;
  elEnonce.textContent = q.enonce;

  // remplir les 4 boutons avec les 4 propositions (i = position du bouton)
  boutonsReponse.forEach((bouton, i) => {
    bouton.textContent = q.propositions[i];
  });
}

// --- Traiter une réponse cliquée ---
function repondre(texteChoisi) {
  const q = questions[indexQuestion];
  if (texteChoisi === q.bonneReponse) {
    score = score + 1;
  }
  indexQuestion = indexQuestion + 1;

  if (indexQuestion < questions.length) {
    afficherQuestion();   // question suivante
  } else {
    terminerQuiz();       // c'était la dernière
  }
}

// --- Fin du quiz ---
function terminerQuiz() {
  elScoreFinal.textContent =
    `Bravo ${pseudo} ! Ton score : ${score} / ${questions.length}`;
  afficherEcran(ecranResultat);
}

// --- Démarrer une partie ---
function demarrerQuiz() {
  pseudo = champPseudo.value;
  if (pseudo === "") {
    pseudo = "Joueur";
  }
  score = 0;
  indexQuestion = 0;
  questions = genererQuestions(10);
  afficherEcran(ecranQuiz);
  afficherQuestion();
}
\end{lstlisting}

\subsection{Étape 5 --- brancher les événements}

C'est ce qui relie les boutons aux fonctions. À placer \textbf{tout en bas} de \code{script.js} :

\begin{lstlisting}[style=js]
// Cliquer sur "Mode Capitales" démarre le quiz
btnCommencer.addEventListener("click", demarrerQuiz);

// Cliquer sur une réponse appelle repondre() avec le texte du bouton
boutonsReponse.forEach((bouton) => {
  bouton.addEventListener("click", (event) => {
    repondre(event.target.textContent);
  });
});

// Cliquer sur "Rejouer" revient à l'accueil
btnRejouer.addEventListener("click", () => {
  afficherEcran(ecranAccueil);
});
\end{lstlisting}

\begin{infobox}[Vérification --- le grand moment]
Lancez Live Server. Saisissez un pseudo, cliquez sur \textbf{Mode Capitales} : l'écran du quiz apparaît avec une question et quatre capitales. Cliquez sur une réponse : la question suivante s'affiche. Après 10 réponses, l'écran Résultat montre votre score. Cliquez sur \textbf{Rejouer} : retour à l'accueil. \textbf{Vous venez de coder votre première application interactive !}
\end{infobox}

% ===================== POUR ALLER PLUS LOIN =====================
\section{Pour aller plus loin (facultatif)}

\begin{objectifbox}[Deux améliorations à tenter]
\begin{enumerate}[leftmargin=*]
  \item \textbf{Compter les erreurs} : ajoutez une variable \code{erreurs} (avec \code{let}), incrémentée dans le \code{else} de \code{repondre}, et affichez-la dans l'écran Résultat.
  \item \textbf{Un petit retour visuel} : ajoutez un élément \code{<p id="feedback"></p>} dans l'écran Quiz, et faites afficher « Bonne réponse ! » ou « Raté » à chaque clic, via \code{textContent}.
\end{enumerate}
\end{objectifbox}

% ===================== AIDE-MÉMOIRE =====================
\section{Aide-mémoire du Lab 7}

\begin{center}
\begin{tabular}{p{5.4cm} p{7.6cm}}
\toprule
\textbf{Besoin} & \textbf{Code} \\
\midrule
Sélectionner un élément & \code{document.querySelector("\#id")} \\
Sélectionner plusieurs & \code{document.querySelectorAll(".classe")} \\
Lire / changer un texte & \code{element.textContent} \\
Lire un champ de saisie & \code{champ.value} \\
Cacher / montrer & \code{element.classList.add/remove("cachee")} \\
Réagir à un clic & \code{element.addEventListener("click", fonction)} \\
Élément cliqué & \code{event.target} \\
Créer / accrocher un élément & \code{createElement(...)} \quad / \quad \code{appendChild(...)} \\
\bottomrule
\end{tabular}
\end{center}

% ===================== GIT =====================
\section{Sauvegarder mon travail (routine Git)}

\begin{retenirbox}[Routine Git de fin de Lab --- tag \texttt{v0.7}]
\begin{lstlisting}[style=html]
git status
git add .
git commit -m "Lab 7 : premier quiz fonctionnel (DOM + evenements)"
git tag v0.7
git push
git push --tags
\end{lstlisting}
Vérifiez sur GitHub que le commit \emph{et} le tag \code{v0.7} apparaissent.
\end{retenirbox}

% ===================== RAPPORT =====================
\begin{attentionbox}[Rapport à produire]
Déposez \code{rapport-lab-7-NOM-Prenom.pdf} dans \code{rapports/}, avec quatre sections :
\begin{enumerate}[leftmargin=*,noitemsep]
  \item \textbf{Ce que j'ai fait} : le quiz à trois écrans, la génération de questions, les événements.
  \item \textbf{Difficultés rencontrées} : vos erreurs (messages copiés). Un quiz a beaucoup de pièces : c'est normal d'avoir des bugs.
  \item \textbf{Comment je les ai résolues}.
  \item \textbf{Lien vers mon repo} avec le tag \code{v0.7}.
\end{enumerate}
\end{attentionbox}

% ===================== RÉSUMÉ =====================
\section{Résumé et transition}

C'est une étape majeure : Jàng Afrig est désormais une \textbf{application interactive}. Vous savez sélectionner des éléments, lire et modifier la page, et réagir aux clics. Vous avez relié HTML, CSS et JavaScript en un tout cohérent.

Mais il y a une limite : le quiz ne connaît que \textbf{dix pays}, écrits à la main dans \code{data.js}. L'Afrique en compte \textbf{54} ! Au \textbf{CM4}, on apprendra comment une page web va chercher de \emph{vraies} données en ligne, en parlant à un service distant (une \emph{API}). Puis, au Lab 8, on remplacera nos dix pays par les 54 pays réels, récupérés automatiquement.

\begin{center}
\textit{Le quiz est jouable. Bientôt, il connaîtra toute l'Afrique.}
\end{center}

\bigskip
\noindent\rule{\textwidth}{0.4pt}
\begin{center}
\small Dr.~El Hadji Bassirou TOURÉ --- UCAD\\
\small Licence 1 DUT1-Informatique --- Fondamentaux du Web --- 2025--2026
\end{center}

\end{document}
